% Source-derived chapter 01: Introduction
% Original source: standard-conjectures-abelian-fourfolds
\chapter{Introduction}
\label{standard-conjectures-abelian-fourfolds-chapter-01}
\source{standard-conjectures-abelian-fourfolds}

\begin{remark}[Source section]
\label{standard-conjectures-abelian-fourfolds-chapter-01-source}
\source{standard-conjectures-abelian-fourfolds}
\source{standard-conjectures-abelian-fourfolds}
This chapter reproduces the corresponding section of the retrieved
LaTeX source, including its definitions, statements, examples, and
proofs. It has not yet been translated into Lean.
\notready
\end{remark}

% The source section title was: Introduction
In this paper we prove that the standard conjecture of Hodge type holds for abelian fourfolds. This is the first  unconditional\footnote{Milne showed that the   Hodge conjecture for complex abelian varieties  implies the standard conjecture of Hodge type for abelian varieties over any field \cite{pola}.}  result on the conjecture since its formulation. Before giving the precise statement and a sketch of the proof, we briefly recall the history of the problem. 



In this introduction $X$ will be a smooth, projective and geometrically connected variety over a base field $k$ of characteristic $p\geq 0$. We denote by  $\mathcal{Z}_{\num}^n(X)_{\Q}$ the space of algebraic cycles of codimension $n$ with $\Q$-coefficients modulo numerical equivalence. This is a finite dimensional {$\Q$-vector} space and its dimension will be denoted by
\begin{equation} \label{eq:intro}
\source{standard-conjectures-abelian-fourfolds}
\rho_n=\dim \mathcal{Z}_{\num}^n(X)_{\Q}. \tag{$*$}
\end{equation}

\subsection*{Brief historical panorama}
The history of the problem starts with the so called Hodge index  theorem.
\begin{theorem}
\source{standard-conjectures-abelian-fourfolds}
\notready\label{introsegre}
\source{standard-conjectures-abelian-fourfolds}
Suppose that $X$ has dimension two. Then the intersection product  
\[\mathcal{Z}^1_{\num}(X)_{\Q}\times \mathcal{Z}_{\num}^1(X)_{\Q} \longrightarrow \Q\]
is of signature $(s_+ ; s_-)=(1; \rho_1 - 1)$.
\end{theorem}
When the characteristic $p$ is zero, the above theorem was proved by Hodge by relating  the intersection product to the cup product in singular cohomology through the cycle class map \cite{Hodge}. An algebraic proof, valid in any characteristic,  was found by Segre \cite{Segre} and Bronowski \cite{Brono}.

Elaborating on an argument of Mattuck--Tate \cite {Mattuck}, Grothendieck  realized that the Lang--Weil estimate for the number of rational points on a smooth and projective curve $C$ over a finite field follows from the Hodge index theorem applied to the surface $X=C\times C$ \cite{Grothseg}. He then proposed a program to show the Weil conjectures for varieties of higher dimension based on a conjectural generalisation of the Hodge index theorem,  known as the standard conjecture of Hodge type. Together with the three other standard conjectures (and the resolution of singularities), it was considered by Grothendieck as the most urgent task in algebraic geometry \cite{GrothSC}. 

This conjecture is connected with other arithmetic contexts such as the conjectural description of the rational points on a Shimura variety over a finite field \cite{LRSC} and the weight-monodromy conjecture \cite{mSaito}. 

The four standard conjectures imply that the category of numerical motives   is semisimple and polarizable just as the category of Hodge structures of smooth projective complex varieties \cite{Rivano}. Surprisingly enough, Jannsen proved semisimplicity   unconditionally and independently of the standard conjectures \cite{Jann}. Polarizability is still open and it is intimately related to the standard conjecture of Hodge type.


 
 As a conclusion of this historical panorama, let us mention that Gillet and Soul\'e have proposed an arithmetic version of the standard conjectures (i.e. over a ring of integers) \cite{soule}. Some results on the arithmetic standard conjecture of Hodge type can be found in \cite{Kunz1,Kunz2,Kunz3,Kresch}.
  
  \subsection*{Main results}
  Let us now recall the formulation of the standard conjecture of Hodge type.
  \begin{conj}
\source{standard-conjectures-abelian-fourfolds}
\notready
  Let   $d$ be the dimension of   a smooth, projective and geometrically connected variety $X$ and fix a hyperplane section $L$ of $X$.  For $n \leq d/2$ define the space of primitive cycles $\mathcal{Z}^{n,\prim}_{\num}(X)_{\Q}$ as
\[\mathcal{Z}^{n,\prim}_{\num}(X)_{\Q} = \{\alpha \in \mathcal{Z}^{n}_{\num}(X)_{\Q} , \hspace{0.2cm} \alpha \cdot L^ {{d-2n+1}}=0 \hspace{0.2cm}  \textrm{in}  \hspace{0.2cm}  \mathcal{Z}^{d-n+1}_{\num}(X)_{\Q}\}\]
and define the pairing  
\[\langle \cdot, \cdot \rangle_{n} : \mathcal{Z}^{n,\prim}_{\num}(X)_{\Q}  \times  \mathcal{Z}^{n,\prim}_{\num}(X)_{\Q}  \longrightarrow \Q\]
 via the intersection product
 \[\alpha, \beta \mapsto  (-1)^n \alpha \cdot \beta \cdot L^{d-2n}.\]
The   standard conjecture of Hodge type predicts that this pairing is positive definite.
\end{conj}
  
 
The evidences for Grothendieck were the case $n=1$ (which can be shown by reducing it to Theorem \ref{introsegre}) and the case     $p=0$. Indeed,  in characteristic zero, one can use the cycle class map to relate the quadratic form $\langle \cdot, \cdot \rangle_{n}$ to the quadratic form given by the cup product on singular cohomology. Then  this kind of positivity statements can be deduced from positivity statements in cohomology, such as the Hodge--Riemann relations in Hodge theory. 



\


 The following is our main theorem. 

\begin{theorem}
\source{standard-conjectures-abelian-fourfolds}
\notready\label{thmintro}
\source{standard-conjectures-abelian-fourfolds}
The standard conjecture of Hodge type holds for abelian fourfolds in positive characteristic.\end{theorem}
It turns out that this statement is equivalent to the following, which is maybe a more direct formulation (see also  Proposition \ref{anypolarization}).
\begin{theorem}
\source{standard-conjectures-abelian-fourfolds}
\notready 
Let $X$ be an abelian fourfold. Then the intersection product 
\[\mathcal{Z}^2_{\num}(X)_{\Q}\times \mathcal{Z}_{\num}^2(X)_{\Q} \longrightarrow \Q\]
is of signature  $ (s_+;s_-) = ( \rho_2 - \rho_1+1;  \rho_1-1)$, with $\rho_n$ as in $($\ref{eq:intro}$)$.
\end{theorem} 
This formulation should   be reminiscent of the Hodge index theorem.


\begin{remark}
\source{standard-conjectures-abelian-fourfolds}
\notready\label{remintro}
\source{standard-conjectures-abelian-fourfolds}
As we explained above, the standard conjecture of Hodge type is known in characteristic zero. Hence Theorem \ref{thmintro} is new only   for those algebraic classes that cannot be lifted to characteristic zero. We discuss the existence of such classes in  Appendix \ref{geometricexamples}.
\end{remark}


By combining Theorem \ref{thmintro} with  a theorem of Clozel  \cite{Clozel} we deduce the following.
\begin{theorem}
\source{standard-conjectures-abelian-fourfolds}
\notready\label{thmclozelintro}
\source{standard-conjectures-abelian-fourfolds}
Let $X$ be an   abelian fourfold. Then numerical equivalence on $X$ coincides with $\ell$-adic homological equivalence on $X$ for infinitely many prime numbers $\ell$. 
\end{theorem}
The fact that homological and numerical equivalence should always coincide is also one   of the four standard conjectures. The two others (namely K\"unneth and Lefschetz) being known for abelian varieties, Theorems \ref{thmintro} and \ref{thmclozelintro} imply that in order to fully understand the standard conjectures for abelian fourfolds what is missing is $\ell$-independency of $\ell$-adic homological equivalence.

\subsection*{Idea of the proof}
The starting point of the proof\footnote{An idea of the strategy is also presented in the report  \cite{Ober}.} of Theorem \ref{thmintro} is a classical product formula from the theory of quadratic forms: let $q$ be a $\Q$-quadratic form, if we know  $q \otimes \Q_\ell$ for all prime numbers $\ell$ then we have information on the signature of $q$, more precisely we know the difference $s_+-s_-$ modulo $8$. When $q$ is the quadratic form  $\langle \cdot, \cdot \rangle_{n}$ as above, then one can hope to compute $q \otimes \Q_\ell$ through the cycle class map.

In characteristic $p=0$ one computes $q \otimes \R$ by directly embedding it into singular cohomology. Our strategy should be thought as a way of circumventing the impossibility (when $p>0$) of embedding $q \otimes \R$ in some Weil cohomology by   instead  embedding all the other completions of   $q$ in  Weil cohomologies.

\



In order to deduce the full signature from the information modulo $8$ one needs to be sure that the rank of $q$ is small. Hence, one cannot apply this strategy to the whole space of algebraic cycles but rather one first decomposes it into smaller quadratic subspaces and then computes $q \otimes \Q_\ell$  for each of those.

To do so, one first   reduces the question to varieties defined over a finite field (this is a classical specialization argument) where abelian varieties are known to always  admit complex multiplication \cite{Tate}.  Then one uses complex multiplication to decompose the space of algebraic cycles into smaller quadratic subspaces. Finally,  it turns out that  when the abelian variety has dimension four then   the subspaces constructed are of rank two (at least those where the problem is not trivial).

\


Technically speaking, we do not compute $q \otimes \Q_\ell$ directly, but rather construct another $\Q$-quadratic form $\tilde{q}$ and try to compare these two. This quadratic form $\tilde{q}$ is constructed as follows. First, the abelian fourfold $X$ admits a lifting $\tilde{X}$ to characteristic zero on which  complex multiplication still acts (Honda--Tate). The action of  complex multiplication  decomposes the singular cohomology  $H_{\sing}^4(\tilde{X},\Q)$ into subspaces. Those are quadratic subspaces endowed with the cup product. To a given factor $q$ of $\mathcal{Z}^2_{\num}(X)_{\Q}$ one associates the factor $\tilde{q}$  of $H_{\sing}^4(\tilde{X},\Q)$ which is (roughly speaking) the same irreducible representation for the action of  complex multiplication.

\

The comparison  $q \otimes \Q_\ell \cong \tilde{q}\otimes \Q_\ell$ holds for all $\ell \neq p$   by smooth proper base change in $\ell$-adic cohomology.  On the other hand one computes $\tilde{q} \otimes \R$ using the Hodge--Riemann relations; it is negative definite if and only if the Hodge types appearing in the Hodge structure $\tilde{q}$ are odd. 

Now the theory of quadratic forms (in particular the product formula on Hilbert symbols) tells us that the positivity of $q \otimes \R$ is equivalent to the following statement. The quadratic forms $q \otimes \Q_p $ and $\tilde{q}\otimes \Q_p$ are   not isomorphic  precisely   if the Hodge types of $\tilde{q}$ are odd. (Note that this equivalence holds because our quadratic spaces have rank two.)

\

This formulation translates the  problem into a question in $p$-adic Hodge theory. To solve it we use the $p$-adic comparison theorem which gives a canonical isomorphism $(q\otimes \Q_p) \otimes B_\crys \cong (\tilde{q}\otimes \Q_p) \otimes B_\crys$ over a   $\Q_p$-algebra $B_{\crys}$. The strategy then consists in  writing explicitly the matrix of the isomorphism with respect to  well chosen bases of the two $\Q_p$-structures (i.e. computing the $p$-adic periods) and then exploit this explicit isomorphism to determine whether the two quadratic forms are  isomorphic also over $\Q_p$. 

The reason for which  a $p$-adic period is in general more computable than a complex one is that it comes equipped with some extra structures (in particular the action of an absolute Frobenius and the de Rham filtration) which sometimes characterize it\footnote{For example, an element of $B_\crys$ which is invariant under the Frobenius and sits in the zeroth step  of the filtration must be an element of $\Q_p$  \cite[Theorem 5.3.7]{Fontp}.}. Moreover, in our particular situation, these extra structures are particularly simple: the absolute Frobenius must act trivially (because $q\otimes \Q_p$ is spanned by algebraic cycles) and the only non-trivial subspace appearing in the de Rham filtration is a line (as these quadratic spaces have rank two).

Once   the matrix $f \in M_{2 \times 2} (B_\crys)$ of $p$-adic periods is explicitly computed we exploit the relation $(q\otimes B_\crys)=  (\tilde{q}\otimes B_\crys) \circ f$. Even though in this equality we only know $f$, this information is enough to determine whether the two quadratic forms  $q\otimes \Q_p$ and $\tilde{q}\otimes \Q_p$ have the same discriminant and represent the same elements of $\Q_p$. Again because we are in rank two these properties control whether  $q\otimes \Q_p$ and $\tilde{q}\otimes \Q_p$ are isomorphic.


\subsection*{Organisation of the paper}
In Section $3$ we recall the standard conjecture of Hodge type and present some first reduction steps, among which the fact that it is enough to work over a finite field. We state Theorem \ref{thmintro} (Theorem \ref{thmscht} in the text) and deduce from it Theorem \ref{thmclozelintro} (Theorem \ref{thmclozel} in the text). In Section $4$ we recall classical results on the motive of an abelian variety. In Section $5$ we show that Theorem \ref{thmscht} holds true for Lefschetz classes on abelian varieties, i.e. those algebraic classes that are linear combinations of intersections of divisors. In Section $6$ we recall that an abelian variety $X$ over a finite field admits complex multiplication and use this to decompose its   motive. We also recall that $X$ together with its complex multiplication lifts to characteristic zero, hence its motivic decomposition lifts too. In Section $7$ we study the interesting subfactors of this decomposition and we call them exotic. By definition, they are those containing algebraic classes which are not Lefschetz classes. The main result of the section says that if $X$ has dimension four then exotic motives have rank two (in the sense that their realizations are cohomology groups of rank two).

From there a somehow independent text starts, in which abelian varieties are not involved anymore: we study motives of rank two, living in mixed characteristic and having algebraic classes in positive characteristic. The main result (Theorem \ref{mainthm}) predicts the signature of the intersection product on those algebraic classes. In Section $8$ we put all pieces together and explain how Theorem \ref{mainthm} implies Theorem \ref{thmscht}. In Section $9$ we recall classical facts   from the theory of quadratic forms and use them to reduce Theorem \ref{mainthm}  to a $p$-adic question. The tool to attack this $p$-adic question is $p$-adic Hodge theory which we recall in Section $10$. In Section $11$ we describe the extra-structures (Frobenius and filtrations) on the   $p$-adic periods appearing and we uniquely characterize them with respect to these extra-structures. This characterization is essential to be able to determine them. We will compute them in  Section $12$. The proof of Theorem \ref{mainthm} is concluded in Section $13$.

Appendix \ref{geometricexamples} contains exemples of exotic motives and in particular of those having algebraic classes that cannot be lifted to characteristic zero.
\subsection*{Acknowledgements}
I would like to thank  Olivier Benoist,  Rutger Noot and Jean-Pierre Wintenberger for useful discussions; Eva Bayer, Giancarlo Lucchini Arteche and Olivier Wittenberg for explanations on  Galois cohomology; Xavier Caruso and Matthew Morrow for explanations on comparison theorems;   Olivier Brinon and Adriano Marmora for their help with $p$-adic periods; Robert Laterveer and Gianluca Pacienza for drawing my attention to the example in Proposition  \ref{propfermat}; Fr\'ed\'eric D\'eglise for pointing out the reference \cite{NiziFred}; Javier Fres\'an, Thomas Kr\"amer, Marco Maculan and Anastasia Prokudina for useful comments on the text.

Finally, I would like to thank the anonymous referee for many useful suggestions and especially for pointing out an argument which allowed to compute the $p$-adic periods in a direct way instead of using \cite[\S 9]{Colmez}. This made the $p$-adic part of the paper more self-contained and as elementary as possible. 

This research was partly supported by the grant ANR--18--CE40--0017 of Agence National de la Recherche.
